\section{Introdução}
\label{sec:introducao}

A classificação automática de alimentos consiste em identificar a categoria
de um alimento a partir de sua imagem. Essa tarefa apresenta desafios devido
à variação de aparência dentro de uma mesma classe e à semelhança visual entre
classes distintas. Diferentes formas de preparo, apresentação, iluminação e
enquadramento podem dificultar o reconhecimento. Neste trabalho, o problema
é tratado como uma tarefa de classificação supervisionada multiclasse, na qual
cada imagem deve ser associada a uma das 101 categorias do conjunto
Food-101~\cite{bossard2014food101}.

A abordagem adotada utiliza transferência de aprendizado com uma
EfficientNet-B0~\cite{tan2019efficientnet} pré-treinada no ImageNet. Sua camada
de classificação é substituída para atender às categorias do Food-101,
preservando os pesos das convoluções. Essa configuração permite investigar
a adaptação de uma rede pré-treinada ao reconhecimento de alimentos sem
ajustar todos os seus parâmetros.

O objetivo do estudo é avaliar o efeito de diferentes configurações de
Batch Normalization e Stochastic Depth sobre o desempenho de classificação.
São comparados quatro experimentos: treinamento apenas do classificador;
atualização adicional das médias e variâncias do Batch Normalization;
treinamento também de seus parâmetros de escala e deslocamento; e ativação
do Stochastic Depth nessa última configuração. A comparação utiliza a mesma
divisão de dados e considera as acurácias top-1 e top-5, além da análise dos
erros por classe.
